% !PRINTED 220
\weberpage{220}

The groups $\mathfrak{A}_{\mu,\mu'}$ are \gesperrt{conjugate subgroups} of $\mathfrak{A}$
(Bd.\ II, § 3); for if $x_{1,0}$ passes by $S$ into $x_{\mu,\mu'}$, and hence $R_{1,0}$ into $R_{\mu,\mu'}$,
then by the group
\[
S^{-1}\mathfrak{A}_{1,0}S
\]
the roots of $R_{\mu,\mu'}$ are permuted only among themselves, and therefore
\weq{6}{S^{-1}\mathfrak{A}_{1,0}S = \mathfrak{A}_{\mu,\mu'}.}

Likewise, if $T$ is a substitution in $\mathfrak{B}$ by which
$x_{1,0}$ passes into $x_{\mu,\mu'}$:
\weq{7}{T^{-1}\mathfrak{B}_{1,0}T = \mathfrak{B}_{\mu,\mu'}.}

The largest common subgroup $\mathfrak{A}_0$ of all the groups $\mathfrak{A}_{\mu,\mu'}$
is determined, according to (4), by
\[
b \equiv 0, \quad c \equiv 0, \quad a \equiv \partial \ (\mathrm{mod}\ n)
\]
and consists therefore of the substitutions:
\[
\begin{pmatrix} a, & 0 \\ 0, & a \end{pmatrix},
\]
in which $a$ is an arbitrary number coprime to $n$. $\mathfrak{A}_0$ is identical
with the largest common subgroup of three of the groups $\mathfrak{A}_{\mu,\mu'}$
that belong to different rows. For $\mathfrak{A}_0$ is then
the largest common subgroup of $\mathfrak{A}_{\mu,\mu'}$, $\mathfrak{A}_{\nu,\nu'}$, $\mathfrak{A}_{\varrho,\varrho'}$ if
the three congruences
\[
\begin{aligned}
c\mu^2 - (\partial - a)\mu\mu' - b\mu'^2 &\equiv 0 \\
c\nu^2 - (\partial - a)\nu\nu' - b\nu'^2 &\equiv 0 \ (\mathrm{mod}\ n) \\
c\varrho^2 - (\partial - a)\varrho\varrho' - b\varrho'^2 &\equiv 0
\end{aligned}
\]
are satisfied \gesperrt{only} under the assumption
\[
\partial \equiv a, \quad b \equiv 0, \quad c \equiv 0 \ (\mathrm{mod}\ n)
\]
This takes place when the determinant
\[
\begin{vmatrix} \mu^2, & \mu\mu', & \mu'^2 \\ \nu^2, & \nu\nu', & \nu'^2 \\ \varrho^2, & \varrho\varrho', & \varrho'^2 \end{vmatrix}
= -(\nu\varrho' - \varrho\nu')(\varrho\mu' - \mu\varrho')(\mu\nu' - \nu\mu')
\]
is relatively prime to $n$.

Likewise, by (5), the largest common subgroup $\mathfrak{B}_0$
of all the $\mathfrak{B}_{\mu,\mu'}$ is the totality of the substitutions:
\weq{8}{\begin{pmatrix} \alpha, & 0 \\ 0, & \alpha \end{pmatrix}}
with the condition:
\weq{9}{\alpha^2 \equiv 1 \ (\mathrm{mod}\ n).}
